\begin{tikzpicture}[
    % --- STYLES ---
    base/.style={draw, thick, align=center, minimum height=1cm, minimum width=3cm, font=\sffamily},
    data/.style={base, trapezium, trapezium left angle=75, trapezium right angle=105, fill=blue!5},
    source/.style={base, ellipse, fill=blue!5},
    process/.style={base, rectangle, rounded corners, fill=gray!10},
    agent/.style={process, fill=purple!15, draw=purple!80!black, thick}, 
    memory/.style={process, shape=cylinder, aspect=0.25, draw, shape border rotate=90, fill=gray!20, minimum height=1.5cm}, 
    script/.style={process, fill=blue!10, draw=blue!80!black, thick, minimum width=4cm}, 
    decision/.style={base, diamond, fill=yellow!15, aspect=2, inner sep=2pt},
    success/.style={base, rectangle, rounded corners, fill=green!15, draw=green!60!black, thick},
    container/.style={draw, dashed, rounded corners, inner sep=20pt, fill=gray!5},
    arrow/.style={->, >=stealth, thick, color=black!80},
    dashed arrow/.style={->, >=stealth, thick, dashed, color=black!80}
]

    % ==========================================
    % RUNTIME TRANSPILATION (Bottom Box)
    % ==========================================
    
    % 1. Core Pipeline
    \node[script] (transpiler) {Deterministic Python\\Transpiler Script};
    \node[source, left=2.5cm of transpiler] (code1) {Code \#1:\\New Pyomo instance};
    \node[source, right=2.5cm of transpiler] (code2) {Code \#2:\\Gurobi-native code};
    
    % 2. Testing Layer
    \node[decision, below=2.5cm of transpiler] (test) {Execute\\Test Suite};
    \node[data] at (code1 |- test) (data) {User-provided data};
    
    % 3. Deployment
    \node[success, below=2cm of test] (deploy) {Use the code};

    % Runtime Arrows
    \draw[arrow] (code1) -- node[above, font=\sffamily\small] {Processed by} (transpiler);
    \draw[arrow] (transpiler) -- node[above, font=\sffamily\small] {Outputs} (code2);
    
    \draw[arrow] (code1.south) -- (test.north west);
    \draw[arrow] (code2.south) -- (test.north east);
    \draw[arrow] (data.east) -- (test.west);
    
    \draw[arrow] (test.south) -- node[right, font=\sffamily\small] {Pass} (deploy.north);

    % ==========================================
    % AGENTIC GENERATION (Top Box)
    % ==========================================
    
    \node[agent, above=3.5cm of transpiler] (agent) {LLM Agent};
    \node[memory] at (code1 |- agent) (repo) {Examples/\\Repository};
    
    % Agentic Arrows
    \draw[dashed arrow] (agent.south) -- node[right, font=\sffamily\small] {Generates / Refactors} (transpiler.north);
    \draw[arrow] (repo.east) -- (agent.west);

    % ==========================================
    % WIDE FEEDBACK LOOPS
    % ==========================================
    
    % Fail Loop: Out the right side of the test suite decision box
    \draw[dashed arrow] (test.east) -- ++(9,0) |- node[pos=0.25, right, font=\sffamily\small, align=center] {Fail:\\Send output\\to agent} (agent.east);
    
    % Save to Repo Loop: Creates a junction node on the "Pass" line to trigger the save action
    % Sweeps 7cm to the left, clear of Code 1 and Data, right into the repository
    \draw[dashed arrow] ($(test.south)!0.5!(deploy.north)$) -| ++(-11,0) |- node[pos=0.25, left, font=\sffamily\small, align=center] {Vetted instance \& code\\saved as dynamic tests\\for the agent} (repo.west);

    % ==========================================
    % BOUNDING BOXES
    % ==========================================
    \begin{scope}[on background layer]
        \node[container, fit=(agent)(repo), label={[font=\sffamily\bfseries, yshift=2mm]above left:Agentic Generation (Stochastic, resource-heavy)}] (phase1) {};
        
        \node[container, fit=(code1)(transpiler)(code2)(data)(test)(deploy), label={[font=\sffamily\bfseries, yshift=2mm]above left:Runtime Transpilation (Deterministic, instantaneous)}] (phase2) {};
    \end{scope}

\end{tikzpicture}